% Part: first-order-logic
% Chapter: completeness
% Section: downward-ls

\documentclass[../../../include/open-logic-section]{subfiles}

\begin{document}

\olfileid{fol}{com}{dls}
\olsection{Löwenheim--Skolem定理}

勒文海姆--斯科伦定理表明，如果一个理论有一个无限模型，那么它也有一个至多可数无穷的模型。
由此可得一个直接推论：一阶逻辑无法表达结构的论域是不可数的——任何在所有不可数结构中都满足的语句或语句集，在某个可数结构中也是满足的。

\begin{thm}
\ollabel{thm:downward-ls} 如果 $\Gamma$ 是一致的，那么它有可数模型，即它在一个论域为有限或可数无穷的结构中可满足。
\end{thm}

\begin{proof}
如果 $\Gamma$ 是一致的，那么完备性定理证明中所构造的结构~$\Struct M$，其论域 $\Domain{M}$ 不大于语言~$\Lang L$ 的项集。因此 $\Struct M$ 是至多可数无穷的。
\end{proof}

\begin{thm}
\ollabel{noidentity-ls} 如果 $\Gamma$ 是不含等词的一阶逻辑语言中的一个一致语句集，那么它有可数无穷模型，即它在一个论域为无穷且可数的结构中可满足。
\end{thm}

\begin{proof}
如果 $\Gamma$ 是一致的并且不含任何出现等词的语句，那么完备性定理证明中所构造的结构~$\Struct M$，其论域 $\Domain{M}$ 与语言~$\Lang L'$ 的项集相同。因为 $\Trm[L']$ 是可数无穷的，所以 $\Struct{M}$ 也是可数无穷的。
\end{proof}

\begin{ex}[斯科伦悖论]
Zermelo（策梅洛）--Fraenkel（弗兰克尔）集合论~$\Log{ZFC}$ 是一个十分强大的框架，几乎所有数学陈述都可以在其中表达，包括关于集合大小的命题。例如，$\Log{ZFC}$ 可以证明实数集~$\Real$ 是不可数的，也可以证明康托尔定理（即任何集合的幂集都大于该集合本身），等等。如果 $\Log{ZFC}$ 是一致的，那么它的所有模型都是无穷的；不仅如此，这些模型都包含一些被该理论断定为不可数的元素（例如，使得 $\Log{ZFC}$ 定理“自然数的幂集存在”为真的那个元素）。根据勒文海姆--斯科伦定理，$\Log{ZFC}$ 也有可数模型——这些模型本身是可数的，却包含着“不可数”集合。
\end{ex}

\end{document}